% Compile with XeLaTeX.
\documentclass[UTF8,fontset=fandol,11pt,a4paper]{ctexart}
\usepackage[margin=24mm]{geometry}
\usepackage{amsmath,amssymb}
\usepackage[unicode,hidelinks]{hyperref}
\hypersetup{pdftitle={Counterexamples on c0 and standard l1},pdfauthor={}}
\usepackage{bookmark}
\usepackage{array}
\setlength{\parindent}{0pt}
\setlength{\parskip}{5pt plus 1pt minus 1pt}
\setlength{\emergencystretch}{2em}
\linespread{1.13}
\allowdisplaybreaks[1]
\providecommand{\tightlist}{\setlength{\itemsep}{0pt}\setlength{\parskip}{0pt}}
\setcounter{secnumdepth}{-1}
\date{}
\begin{document}
\begin{center}
{\Large\bfseries $c_0$ 与标准 $\ell^1$ 上的反例}\par\medskip
\end{center}

本文给出指数尾曲线构造的 $c_0$ 反例，以及有限支撑曲线构造的 $c_0$ 反例和它在标准 $\ell^1$ 上的提升。各例均由两个极大单调算子组成，满足原范数拓扑下的内部条件，而其和非极大单调。

\section*{定义与符号}

令 $\mathbb N=\{1,2,\ldots\}$，$\mathbb N_0=\{0,1,\ldots\}$。对实 Banach 空间 $Y$，以 $Y^*$ 表示其连续对偶，并记 $\langle x,a\rangle=a(x)$。对算子 $T:Y\rightrightarrows Y^*$，定义 $\operatorname{gra}T=\{(x,a):a\in Tx\}$ 和 $\operatorname{dom}T=\{x:Tx\ne\varnothing\}$。算子的逐点和定义为 $(T+R)x=\{a+b:a\in Tx,\ b\in Rx\}$。

图 $G\subset Y\times Y^*$ 单调是指其任意两点满足 $\langle x-y,a-b\rangle\ge0$。若它单调且在 $Y\times Y^*$ 中没有真单调扩张，则称其极大单调。定义单调极性集
\[
G^\mu=\{(z,p)\in Y\times Y^*:
\langle z-x,p-a\rangle\ge0\quad\forall(x,a)\in G\}.
\]
单调图 $G$ 极大单调当且仅当 $G=G^\mu$：$G^\mu\setminus G$ 中的点可以加入而保持单调性，任一单调扩张中的点也必须属于 $G^\mu$。特别地，$(0,0)\in G^\mu$ 当且仅当 $\langle x,a\rangle\ge0$ 对所有 $(x,a)\in G$ 成立。

当 $\operatorname{gra}T\ne\varnothing$ 且 $0\notin\operatorname{dom}T$ 时，记
\[
\mathcal V(T)=\sup_{(x,a)\in\operatorname{gra}T}
\frac{\max\{0,-\langle x,a\rangle\}}{\|x\|}.
\]
因此，对每个 $C\ge0$，$\mathcal V(T)\le C$ 当且仅当 $\langle x,a\rangle\ge-C\|x\|$ 对全部图点成立。对任意非空图，还记
\[
\begin{aligned}
F_T(x,p)&=\sup_{(y,a)\in\operatorname{gra}T}
\bigl(\langle x,a\rangle+\langle y,p\rangle-\langle y,a\rangle\bigr),\\
\Phi_T(x,p)&=F_T(x,p)-\langle x,p\rangle\\
&=\sup_{(y,a)\in\operatorname{gra}T}\langle x-y,a-p\rangle.
\end{aligned}
\]
对两个极大单调算子 $T,R$，定义
\[
\Delta_{T,R}(x,p)=\inf_{q\in Y^*}
\bigl[\Phi_T(x,p-q)+\Phi_R(x,q)\bigr].
\]
下文 $g\otimes g$ 表示算子 $x\mapsto\langle x,g\rangle g$，$q_+=\max\{q,0\}$，$c_{00}$ 表示有限支撑序列的集合。所有极大性均在完整连续对偶对中，所有内部均相对于原范数拓扑。

\part*{Counterexample on $c_0$}

在通常的实 Banach 空间 $c_0$ 及其完整对偶 $\ell^1$ 上，构造极大单调算子 $A$ 和非零秩一正算子 $P=g\otimes g$，使 \[
\boxed{
(0,0)\in
\bigl(\operatorname{gra}(A+P)\bigr)^\mu
\setminus\operatorname{gra}(A+P),
}
\] 而且原点对\textbf{整个和图}具有统一严格余量： \[
\boxed{
\inf_{(x,a)\in\operatorname{gra}A}
\bigl(\langle x,a\rangle+g(x)^2\bigr)
=
1-\frac1{16}\sum_{n=1}^{\infty}\frac1{n^2}
>\frac78.
}
\]

\section{一、完整空间和块算子}\label{ux4e00ux5b8cux6574ux7a7aux95f4ux548cux5757ux7b97ux5b50}

取 \[
I=\mathbb N_0\times\mathbb N,\qquad
X=c_0(I),\qquad X^*=\ell^1(I).
\] 这里的 $c_0(I)$ 是\textbf{全指标集上的通常 $c_0$}，不是仅要求每个块分别趋零的较大空间。因为 $I$ 可数，$X$ 与标准 $c_0$ 等距。

对 $a\in\ell^1(I)$，定义 \[
(ma)_b=\sum_{j\ge1}a_{b,j},
\qquad
(Ea)_{b,j}=a_{b,j}+2\sum_{k>j}a_{b,k}.
\tag{1}
\] 于是 \[
m:\ell^1(I)\to\ell^1(\mathbb N_0),
\qquad
E:\ell^1(I)\to c_0(I),
\] 且 \[
\|ma\|_1\le\|a\|_1,\qquad
\|Ea\|_\infty\le2\|a\|_1.
\]

其中 $E$ 的全局 $c_0$ 性质可以直接检查：只有有限多个块的 $\ell^1$ 范数超过任一给定阈值；每个这样的块内，尾和又趋于零。

绝对收敛的求和给出 \[
\boxed{
\langle Ea,b\rangle+\langle Eb,a\rangle
=2\sum_{n\ge0}(ma)_n(mb)_n,
\qquad
\langle Ea,a\rangle=\|ma\|_2^2.
}
\tag{2}
\] 这里的 $\ell^2$ 只用于计算块质量的能量；原始空间及完整对偶始终是 $c_0(I)\times\ell^1(I)$。

固定 \[
g=e_{0,1}-e_{0,2}\in X^*,
\qquad h=Eg\in X.
\] 具体而言，$h$ 的前两个第零块坐标都是 $-1$，其余坐标为零。因此 \[
mg=0,\qquad \langle h,g\rangle=0.
\] 令 \[
r(a)=\langle h,a\rangle=-a_{0,1}-a_{0,2},
\qquad
La=-Ea+r(a)h.
\tag{3}
\]

由式 (2)，得到三个后面要使用的恒等式： \[
\boxed{
\langle La,a\rangle=r(a)^2-\|ma\|_2^2,
}
\tag{4}
\] \[
\boxed{
\langle La-Lb,a-b\rangle
=(r(a)-r(b))^2-\|ma-mb\|_2^2,
}
\tag{5}
\] \[
\boxed{\langle La,g\rangle=r(a).}
\tag{6}
\]

\textbf{不声称 $L$ 在整个 $\ell^1(I)$ 上单调。} 接下来将明确指定其非线性限制域。

\section{二、实际谱曲线与不可成为实际质量的端点}\label{ux4e8cux5b9eux9645ux8c31ux66f2ux7ebfux4e0eux4e0dux53efux6210ux4e3aux5b9eux9645ux8d28ux91cfux7684ux7aefux70b9}

对 $s>0$，定义 \[
\omega_s
=
\left(\frac{e^{-ns^2}}{4n}\right)_{n\ge1}
\in\ell^1(\mathbb N).
\tag{7}
\] 每个实际参数的可求和性是明确的： \[
\|\omega_s\|_1
=-\frac14\log(1-e^{-s^2})<\infty.
\]

另一方面， \[
\omega_s\longrightarrow
\omega_0:=\left(\frac1{4n}\right)_{n\ge1}
\quad\text{in }\ell^2
\quad(s\downarrow0),
\] 但是 \[
\boxed{\omega_0\in\ell^2\setminus\ell^1.}
\tag{8}
\] 记 \[
\sigma=\|\omega_0\|_2^2
=\frac1{16}\sum_{n\ge1}\frac1{n^2}
=\frac{\pi^2}{96}<\frac18,
\qquad
\delta=1-\sigma>\frac78.
\tag{9}
\]

曲线具有统一的 $\ell^2$ Lipschitz 估计。逐坐标求导： \[
\omega_s'(n)=-\frac{s}{2}e^{-ns^2},
\] 所以 \[
\sum_{n\ge1}|\omega_s'(n)|^2
=
\frac{s^2}{4(e^{2s^2}-1)}
\le\frac18.
\] 先对有限截断使用积分 Cauchy-Schwarz，再令截断长度趋于无穷，得到 \[
\boxed{
\|\omega_s-\omega_t\|_2^2\le\frac18(s-t)^2.
}
\tag{10}
\]

现在令 \[
U=(-\infty,-1)\cup(1,\infty),
\qquad
F(t)=\bigl(\operatorname{sgn}t,\omega_{|t|-1}\bigr)
\in\ell^1(\mathbb N_0).
\tag{11}
\] $F$ 在 $\ell^2$ 距离下为 $1$-Lipschitz。同号时直接由式 (10) 得到；异号时写成 $t=1+s$、$u=-1-v$，其中 $s,v>0$，则 \[
\begin{aligned}
\|F(t)-F(u)\|_2^2
&=4+\|\omega_s-\omega_v\|_2^2\\
&\le4+\frac18(s-v)^2\\
&\le(2+s+v)^2=(t-u)^2.
\end{aligned}
\tag{12}
\]

两侧的能量极限是 \[
F_+=(1,\omega_0),\qquad F_-=(-1,\omega_0).
\tag{13}
\] \textbf{二者都不属于实际质量空间 $\ell^1(\mathbb N_0)$。后面只用它们取必要不等式的极限，不将它们插入原图。}

\section{三、明确指定整张原图，并证明单调性}\label{ux4e09ux660eux786eux6307ux5b9aux6574ux5f20ux539fux56feux5e76ux8bc1ux660eux5355ux8c03ux6027}

定义 \[
D=
\left\{
a\in\ell^1(I):
r(a)\in U,\quad ma=F(r(a))
\right\}.
\tag{14}
\] 这里包括满足等式的\textbf{全部} $\ell^1$ 数组。

令 \[
\boxed{
\operatorname{gra}A=\{(La,a):a\in D\}.
}
\tag{15}
\]

每个 $t\in U$ 都有一个明确的实际参数 $a^t\in D$： \[
a^t_{0,1}=-t,\qquad
a^t_{0,2}=0,\qquad
a^t_{0,3}=t+\operatorname{sgn}t,
\] \[
a^t_{n,1}
=\frac{e^{-n(|t|-1)^2}}{4n}
\quad(n\ge1),
\] 其余坐标为零。

于是 \[
r(a^t)=t,\qquad ma^t=F(t),\qquad a^t\in\ell^1(I).
\] 这既证明 $D\ne\varnothing$，也保证以后对任意 $t\in U$ 的测试都有\textbf{原图实际行}。

对全部 $a,b\in D$，由式 (5)、(12)， \[
\begin{aligned}
\langle La-Lb,a-b\rangle
&=(r(a)-r(b))^2
-\|F(r(a))-F(r(b))\|_2^2\\
&\ge0.
\end{aligned}
\tag{16}
\] 因此 $A$ 单调。

\section{\texorpdfstring{四、核心：完整 $c_0\times\ell^1$ 对偶对中的极大性}{四、核心：完整 c\_0\textbackslash times\textbackslash ell\^{}1 对偶对中的极大性}}\label{ux56dbux6838ux5fc3ux5b8cux6574-c_0timesell1-ux5bf9ux5076ux5bf9ux4e2dux7684ux6781ux5927ux6027}

任取实际查询 \[
(x,p)\in c_0(I)\times\ell^1(I)
\] 满足 \[
\langle x-La,p-a\rangle\ge0
\qquad\forall a\in D.
\tag{17}
\] 下面证明它必定属于式 (15) 的原图。

\subsection{4.1 全部实际核直线压出查询坐标}\label{ux5168ux90e8ux5b9eux9645ux6838ux76f4ux7ebfux538bux51faux67e5ux8be2ux5750ux6807}

令 \[
K=\{k\in\ell^1(I):mk=0,\ r(k)=0\}.
\] 固定任意 $a\in D$。对所有 $k\in K$、$\theta\in\mathbb R$， \[
a+\theta k\in D.
\] 因此可以将这条完整实际直线代入式 (17)。

因为 $Lk=-Ek$ 且 $\langle Lk,k\rangle=0$，直接展开得到 \[
\boxed{
\langle x-L(a+\theta k),p-a-\theta k\rangle
=
\langle x-La,p-a\rangle
-\theta\langle x+Ep,k\rangle.
}
\tag{18}
\] 对正负全部 $\theta$ 非负，便迫使 \[
\langle x+Ep,k\rangle=0\qquad(k\in K).
\tag{19}
\]

设 $w=x+Ep\in c_0(I)$。只需使用有限支撑的实际核方向：

对于每个 $b\ge1$，坐标差 $e_{b,j}-e_{b,k}\in K$，故 $w$ 在该无限块内恒定；由 $c_0$ 性质，该块恒为零。第零块从第三坐标起同理全为零。再用 $g=e_{0,1}-e_{0,2}\in K$，得到前两个坐标相等。

因此存在 $v\in\mathbb R$，使 \[
\boxed{x+Ep=vh,\qquad x=-Ep+vh.}
\tag{20}
\]

\subsection{4.2 全图 polar 条件化为精确平方不等式}\label{ux5168ux56fe-polar-ux6761ux4ef6ux5316ux4e3aux7cbeux786eux5e73ux65b9ux4e0dux7b49ux5f0f}

记 \[
r_p=r(p),\qquad m_p=mp\in\ell^1\subset\ell^2.
\] 由式 (2)、(20)，对每条实际行有 \[
\boxed{
\langle x-La,p-a\rangle
=
(r(a)-v)(r(a)-r_p)
-\|F(r(a))-m_p\|_2^2.
}
\tag{21}
\]

每个 $t\in U$ 都有实际 $a^t$，故式 (17) 等价于 \[
\|F(t)-m_p\|_2^2\le(t-v)(t-r_p)
\qquad(t\in U).
\] 令 \[
\tau=\frac{v+r_p}{2},\qquad
\nu=\frac{v-r_p}{2}.
\] 得到 \[
\boxed{
\|F(t)-m_p\|_2^2+\nu^2
\le(t-\tau)^2
\qquad(t\in U).
}
\tag{22}
\]

\subsection{4.3 外侧查询直接返回实际行}\label{ux5916ux4fa7ux67e5ux8be2ux76f4ux63a5ux8fd4ux56deux5b9eux9645ux884c}

若 $|\tau|>1$，则 $\tau\in U$，可以在式 (22) 中代入\textbf{实际参数} $t=\tau$。于是 \[
m_p=F(\tau),\qquad \nu=0.
\] 从而 \[
v=r_p=\tau,\qquad p\in D,\qquad
x=-Ep+r_ph=Lp.
\] 因此 \[
(x,p)\in\operatorname{gra}A.
\]

\subsection{4.4 中间查询被完整对偶排除}\label{ux4e2dux95f4ux67e5ux8be2ux88abux5b8cux6574ux5bf9ux5076ux6392ux9664}

若 $|\tau|\le1$，分别令实际参数 $t\downarrow1$ 和 $t\uparrow-1$。由 $\ell^2$ 收敛，式 (22) 要求 \[
\|F_+-m_p\|_2^2+\nu^2\le(1-\tau)^2,
\] \[
\|F_--m_p\|_2^2+\nu^2\le(1+\tau)^2.
\tag{23}
\]

写 \[
m_p=(b,z),\qquad z\in\ell^1(\mathbb N),
\] 并令 \[
R=\|z-\omega_0\|_2^2+\nu^2\ge0.
\] 那么式 (23) 是 \[
(b-1)^2+R\le(1-\tau)^2,
\qquad
(b+1)^2+R\le(1+\tau)^2.
\tag{24}
\]

由于 $-1\le\tau\le1$，第一式给出 $b\ge\tau$，第二式给出 $b\le\tau$。所以 $b=\tau$，代回即得 \[
R=0,\qquad z=\omega_0.
\] 但这与 \[
z\in\ell^1,\qquad\omega_0\notin\ell^1
\] 矛盾。

\textbf{这里不是用一个非实际端点充当返回行；恰恰相反，是证明任何中间的原空间 polar 查询都必须携带一个不可能的实际 $\ell^1$ 质量。}

两种情形穷尽全部实际查询，因此 \[
\boxed{
(\operatorname{gra}A)^\mu=\operatorname{gra}A.
}
\tag{25}
\] 这就是原始完整 Banach 对偶对中的极大性。

\section{五、同一张图上的秩一和失效}\label{ux4e94ux540cux4e00ux5f20ux56feux4e0aux7684ux79e9ux4e00ux548cux5931ux6548}

定义 \[
\boxed{Px=\langle x,g\rangle g.}
\tag{26}
\] $P$ 非零、秩一、有界、全定义，且 \[
\langle x-y,Px-Py\rangle
=\langle x-y,g\rangle^2\ge0.
\]

它自身的极大性也直接成立：对一个 $P$-polar 查询 $(z,z^*)$，测试 $(z+tu,P(z+tu))$，得 \[
-t\langle u,z^*-Pz\rangle
+t^2\langle u,Pu\rangle\ge0.
\] 令 $t$ 从正负两侧趋零，得到 $z^*=Pz$。

现在，对全部 $a\in D$，式 (6) 给出 \[
g(La)=r(a),\qquad |r(a)|>1.
\] 因此 \[
\boxed{0\notin\operatorname{dom}A.}
\tag{27}
\]

而整个和图的原点括号为 \[
\begin{aligned}
\langle La,a+P(La)\rangle
&=2r(a)^2-\|F(r(a))\|_2^2\\
&=2r(a)^2-1
-\frac1{16}\sum_{n\ge1}
\frac{e^{-2n(|r(a)|-1)^2}}{n^2}\\
&>\delta=1-\sigma>\frac78.
\end{aligned}
\tag{28}
\]

所以 \[
\boxed{
(0,0)\in
(\operatorname{gra}(A+P))^\mu
\setminus\operatorname{gra}(A+P).
}
\tag{29}
\]
由于 $A+P$ 单调，加入 $(0,0)$ 后得到其真单调扩张，故 $A+P$ 非极大单调。

资格条件完整成立： \[
\boxed{
\operatorname{dom}A
\cap\operatorname{int}\operatorname{dom}P
=
\operatorname{dom}A\ne\varnothing.
}
\tag{30}
\] 这给出的四元组是 \[
\boxed{\bigl(c_0(I),A,g\otimes g,(0,0)\bigr).}
\]

式 (28) 的下确界精确等于 $\delta$：用已经写出的实际 $a^t$，令 $t\downarrow1$ 即可。没有要求一个非实际端点达到下确界。

\section{六、逆纤维、缩放与反馈映射}\label{ux516dux4e0eux5f53ux524d-sinf-ux548cux53cdux9988ux95eeux9898ux7684ux7cbeux786eux8fdeux63a5}

逆纤维单点性直接包含在定义里： \[
A^{-1}(a)=
\begin{cases}
\{La\},&a\in D,\\
\varnothing,&a\notin D.
\end{cases}
\] 因此每个逆纤维至多一点。

令 $c=\delta^{-1/2}$，定义 \[
\widehat A(x)=c\,A(x/c).
\] 它的图是 \[
\{(cLa,ca):a\in D\}.
\] 全部差括号乘以 $c^2$，所以仍极大单调，并保持单点逆纤维。同时 \[
\boxed{
\inf_{(x,a)\in\operatorname{gra}\widehat A}
\bigl(\langle x,a\rangle+g(x)^2\bigr)=1.
}
\tag{31}
\] 因此该缩放保留极大单调性及逆纤维单点性，并将上述下确界归一化为 $1$。

未缩放时，$B=A^{-1}=L|_D$ 的反馈也完全明确： \[
\Phi_{\mathrm{fb}}(a)=a+g(Ba)g=a+r(a)g.
\] 由于 $mg=0$、$r(g)=0$， \[
\boxed{\Phi_{\mathrm{fb}}(D)=D,\qquad \Phi_{\mathrm{fb}}^{-1}(a)=a-r(a)g,\qquad0\notin D.}
\tag{32}
\]

每条已经遇到 $D$ 的 $g$-直线，实际上都完整地包含在 $D$： \[
a_0+\mathbb Rg\subset D\qquad(a_0\in D),
\] \[
B(a_0+tg)=Ba_0-th,\qquad
g(B(a_0+tg))=r(a_0).
\tag{33}
\] 但目标直线 $\mathbb Rg$ 完全避开 $D$。

由自配对公式和实际参数 $t\downarrow1$，还可计算
\[
\begin{aligned}
F_A(0,0)
&=\sup_{|t|>1}\bigl(1+\|\omega_{|t|-1}\|_2^2-t^2\bigr)=\sigma,\\
F_{A+P}(0,0)
&=\sup_{|t|>1}\bigl(1+\|\omega_{|t|-1}\|_2^2-2t^2\bigr)=\sigma-1.
\end{aligned}
\]
因此
\[
F_A(0,0)-F_{A+P}(0,0)=1.
\]

\clearpage
\part*{Finite-support version and counterexample on standard $\ell^1$}

下面将指数尾替换为每个实际参数下均有限支撑的曲线，重新证明所得算子的极大性，再通过一个有界线性满射构造标准 $\ell^1$ 上的反例。

\section{一、要证明的反例结论}\label{ux4e00ux8981ux8bc1ux660eux7684ux53cdux4f8bux7ed3ux8bba}

将构造 \[
T:\ell^1(\mathbb N)\rightrightarrows\ell^\infty(\mathbb N),
\] 满足 \[
\boxed{
\begin{gathered}
T\text{ 在完整 }\ell^1\times\ell^\infty\text{ 中极大单调},\\
T(0)=\varnothing,\\
\langle u,b\rangle\ge-\frac14\|u\|_1
\qquad\forall(u,b)\in\operatorname{gra}T.
\end{gathered}}
\]

还将构造一个非零正秩一、全定义、极大单调算子 \[
Bu=\langle u,f\rangle f
\] 使 \[
\boxed{
(0,0)\in
\bigl(\operatorname{gra}(T+B)\bigr)^\mu
\setminus\operatorname{gra}(T+B).
}
\]

由于 $B$ 的定义域为整个 $\ell^1$，只要 $\operatorname{dom}T\ne\varnothing$，便满足 $\operatorname{dom}T\cap\operatorname{int}\operatorname{dom}B\ne\varnothing$。上述性质因此给出原内部条件下一般求和命题的反例，并同时给出 $\mathcal V(T)<\infty$ 而 $T(0)=\varnothing$ 的极大单调算子。

\section{\texorpdfstring{二、先在 $c_0$ 的完整对偶对中构造算子}{二、先在 c\_0 的完整对偶对中构造算子}}\label{ux4e8cux5148ux5728-c_0-ux7684ux5b8cux6574ux5bf9ux5076ux5bf9ux4e2dux6784ux9020ux7b97ux5b50}

\subsection{1. 两层可数指标与线性配对}\label{ux4e24ux5c42ux53efux6570ux6307ux6807ux4e0eux7ebfux6027ux914dux5bf9}

令 \[
I=\mathbb N_0\times\mathbb N,\qquad
X=c_0(I),\qquad X^*=\ell^1(I).
\]

这里的 $c_0(I)$ 是\textbf{整个指标集上的通常 $c_0$}，不是只要求每一块分别趋零。

对 $a=(a_{b,j})\in\ell^1(I)$，定义 \[
(ma)_b=\sum_{j\ge1}a_{b,j},
\qquad
(Ea)_{b,j}=a_{b,j}+2\sum_{k>j}a_{b,k}.
\]

有 \[
\|ma\|_1\le\|a\|_1,\qquad
\|Ea\|_\infty\le2\|a\|_1.
\]

而且 $Ea\in c_0(I)$：对于任何正阈值，只有有限多个块的 $\ell^1$ 质量超过该阈值；每个这样的块内，尾和趋零。其余块的全部坐标则由两倍块质量统一控制。

绝对收敛的求和给出 \[
\langle Ea,b\rangle+\langle Eb,a\rangle
=2\langle ma,mb\rangle_{\ell^2},
\] 特别是 \[
\langle Ea,a\rangle=\|ma\|_2^2.
\tag{1}
\]

取 \[
g=e_{0,1}-e_{0,2},\qquad h=Eg.
\] 则 $h$ 只有第零块的前两个坐标非零，二者均为 $-1$，并且 \[
mg=0,\qquad \langle h,g\rangle=0.
\]

定义 \[
r(a)=\langle h,a\rangle=-a_{0,1}-a_{0,2},
\qquad
La=-Ea+r(a)h.
\]

由式（1）， \[
\boxed{
\langle La-Lb,a-b\rangle
=(r(a)-r(b))^2-\|ma-mb\|_2^2,
}
\tag{2}
\] 以及 \[
\boxed{
\langle La,a\rangle=r(a)^2-\|ma\|_2^2,
\qquad
\langle La,g\rangle=r(a).
}
\tag{3}
\]

这些线性映射与第一部分相同，下面采用不同的参数曲线。

\subsection{2. 用有限支撑尾部实现非实际端点}\label{ux7528ux6709ux9650ux652fux6491ux5c3eux90e8ux5b9eux73b0ux975eux5b9eux9645ux7aefux70b9}

对 $s>0$，定义 \[
\omega_s(n)=\frac{(1-sn^{1/4})_+}{4n},
\qquad n\ge1.
\tag{4}
\]

因为只有 $n<s^{-4}$ 时可能非零， \[
\omega_s\in c_{00}\subset\ell^1
\qquad(s>0).
\]

但当 $s\downarrow0$ 时， \[
\omega_s\longrightarrow
\omega_0:=\left(\frac1{4n}\right)_{n\ge1}
\quad\text{在 }\ell^2\text{ 范数中},
\] 而 \[
\omega_0\in\ell^2\setminus\ell^1.
\tag{5}
\]

记 \[
\sigma:=\|\omega_0\|_2^2
=\frac1{16}\sum_{n\ge1}\frac1{n^2}
<\frac18.
\]

正部函数是 $1$-Lipschitz，故 \[
|\omega_s(n)-\omega_t(n)|
\le\frac{|s-t|}{4n^{3/4}}.
\] 从而 \[
\|\omega_s-\omega_t\|_2^2
\le\frac1{16}
\left(\sum_{n\ge1}n^{-3/2}\right)|s-t|^2
\le\frac3{16}|s-t|^2.
\tag{6}
\]

这里 $\sum_{n\ge1}n^{-3/2}<1+\int_1^\infty x^{-3/2}\,dx=3$。

令 \[
J=(-\infty,-1)\cup(1,\infty),
\] \[
F(t)=\bigl(\operatorname{sgn}t,\omega_{|t|-1}\bigr)
\in\ell^1(\mathbb N_0),
\qquad t\in J.
\tag{7}
\]

这个 $F$ 满足 \[
\|F(t)-F(u)\|_2\le|t-u|.
\tag{8}
\]

同号时由式（6）直接得到。异号时写成 $t=1+s$、$u=-1-v$，则 \[
\begin{aligned}
\|F(t)-F(u)\|_2^2
&=4+\|\omega_s-\omega_v\|_2^2\\
&\le4+\frac3{16}(s-v)^2\\
&\le(2+s+v)^2.
\end{aligned}
\]

现在定义\textbf{完整参数域} \[
D=\{a\in\ell^1(I):r(a)\in J,\ ma=F(r(a))\},
\tag{9}
\] 并定义整张图 \[
\boxed{
\operatorname{gra}A=\{(La,a):a\in D\}.
}
\tag{10}
\]

这个域不是假想的。每个 $t\in J$ 都有一个明确的有限支撑参数： \[
a^t_{0,1}=-t,\quad
a^t_{0,2}=0,\quad
a^t_{0,3}=t+\operatorname{sgn}t,
\] \[
a^t_{n,1}=\omega_{|t|-1}(n)\quad(n\ge1),
\] 其余坐标为零。它满足 \[
r(a^t)=t,\qquad ma^t=F(t).
\]

而对任意 $a,b\in D$，式（2）、（8）给出 \[
\langle La-Lb,a-b\rangle
=(r(a)-r(b))^2
-\|F(r(a))-F(r(b))\|_2^2
\ge0.
\] 所以 $A$ 非空且单调。

\section{三、关键：证明它是完整极大单调图}\label{ux4e09ux5173ux952eux8bc1ux660eux5b83ux662fux5b8cux6574ux6781ux5927ux5355ux8c03ux56fe}

这一部分不能由有限截断、数值实验或``选一个极大扩张''替代。

任取 \[
(x,p)\in c_0(I)\times\ell^1(I)
\] 满足 \[
\langle x-La,p-a\rangle\ge0
\qquad\forall a\in D.
\tag{11}
\]

需要证明 $(x,p)$ 已经属于式（10）。

\subsection{1. 全部实际核直线强迫查询的形式}\label{ux5168ux90e8ux5b9eux9645ux6838ux76f4ux7ebfux5f3aux8febux67e5ux8be2ux7684ux5f62ux5f0f}

令 \[
K=\{k\in\ell^1(I):mk=0,\ r(k)=0\}.
\]

对任意 $a\in D$、$k\in K$、$\theta\in\mathbb R$，都有 \[
a+\theta k\in D.
\]

将这条\textbf{完整实际直线}代入式（11）。由式（1）展开得 \[
\begin{aligned}
&\langle x-L(a+\theta k),p-a-\theta k\rangle\\
&\qquad=
\langle x-La,p-a\rangle
-\theta\langle x+Ep,k\rangle.
\end{aligned}
\tag{12}
\]

因为这对全部正负 $\theta$ 非负，必有 \[
\langle x+Ep,k\rangle=0
\qquad(k\in K).
\tag{13}
\]

设 $w=x+Ep\in c_0(I)$。

对每个 $b\ge1$，所有坐标差 \[
e_{b,j}-e_{b,l}\in K,
\] 因此 $w$ 在该无限块内恒定。由于 $w\in c_0(I)$，这个常数只能为零。

第零块从第三个坐标起同理全为零。再用 \[
e_{0,1}-e_{0,2}=g\in K
\] 得到前两个坐标相同。

所以存在 $v\in\mathbb R$，使 \[
\boxed{x+Ep=vh,\qquad x=-Ep+vh.}
\tag{14}
\]

\subsection{2. 全图极性化为精确平方不等式}\label{ux5168ux56feux6781ux6027ux5316ux4e3aux7cbeux786eux5e73ux65b9ux4e0dux7b49ux5f0f}

记 \[
r_p=r(p),\qquad m_p=mp\in\ell^1\subset\ell^2.
\]

对任意实际参数 $a\in D$，令 $t=r(a)$。由式（14）， \[
\begin{aligned}
\langle x-La,p-a\rangle
&=-\|m_p-F(t)\|_2^2+(v-t)(r_p-t)\\
&=(t-v)(t-r_p)-\|F(t)-m_p\|_2^2.
\end{aligned}
\tag{15}
\]

每个 $t\in J$ 都有实际参数 $a^t$ 可测试。因此，令 \[
\tau=\frac{v+r_p}{2},
\qquad
\nu=\frac{v-r_p}{2},
\] 式（11）强迫 \[
\boxed{
\|F(t)-m_p\|_2^2+\nu^2
\le(t-\tau)^2
\qquad\forall t\in J.
}
\tag{16}
\]

现在只有两种情况。

\textbf{若 $|\tau|>1$}，则 $\tau\in J$，直接代入实际参数 $t=\tau$，得到 \[
m_p=F(\tau),\qquad \nu=0.
\] 所以 \[
v=r_p=\tau,\qquad p\in D,\qquad x=Lp.
\] 即查询已经在原图内。

\textbf{若 $|\tau|\le1$}，分别让实际参数 $t\downarrow1$、$t\uparrow-1$。由式（5），得到 \[
\|(1,\omega_0)-m_p\|_2^2+\nu^2
\le(1-\tau)^2,
\] \[
\|(-1,\omega_0)-m_p\|_2^2+\nu^2
\le(1+\tau)^2.
\tag{17}
\]

写 \[
m_p=(b,z),\qquad z\in\ell^1,
\] 并令 \[
R=\|z-\omega_0\|_2^2+\nu^2\ge0.
\] 则 \[
(b-1)^2+R\le(1-\tau)^2,
\qquad
(b+1)^2+R\le(1+\tau)^2.
\]

因为 $-1\le\tau\le1$，第一式给出 $b\ge\tau$，第二式给出 $b\le\tau$。因此 \[
b=\tau,\qquad R=0.
\] 特别是 \[
z=\omega_0,
\] 这与 \[
z\in\ell^1,\qquad\omega_0\notin\ell^1
\] 矛盾。

\textbf{这里没有将非实际端点插入图中；恰恰是在证明，任何中间的实际 polar 查询都要求一个不可能的 $\ell^1$ 质量。}

因此， \[
\boxed{
(\operatorname{gra}A)^\mu=\operatorname{gra}A.
}
\tag{18}
\]

这证明了原始完整 $c_0(I)\times\ell^1(I)$ 对偶对中的极大性。

\section{四、有限斜率与缺失原点同时成立}\label{ux56dbux6709ux9650ux659cux7387ux4e0eux7f3aux5931ux539fux70b9ux540cux65f6ux6210ux7acb}

对任何 $a\in D$，令 $t=r(a)$。由式（3）， \[
|\langle La,g\rangle|=|t|>1.
\] 由于 $\|g\|_1=2$， \[
\|La\|_\infty>\frac12.
\tag{19}
\]

所以 \[
0\notin\operatorname{dom}A.
\]

另一方面， \[
\begin{aligned}
\langle La,a\rangle
&=t^2-\|F(t)\|_2^2\\
&=t^2-1-\|\omega_{|t|-1}\|_2^2\\
&>-\sigma.
\end{aligned}
\]

结合式（19），得到对\textbf{全部实际图行} \[
\boxed{
\langle La,a\rangle
\ge-2\sigma\|La\|_\infty,
\qquad 2\sigma<\frac14.
}
\tag{20}
\]

至此，已经不是``有限斜率种子可能怎样极大化''，而是整张已证明极大的图本身具有有限斜率，且缺失原点。

\section{\texorpdfstring{五、把这个完整对象提升到标准 $\ell^1$}{五、把这个完整对象提升到标准 \textbackslash ell\^{}1}}\label{ux4e94ux628aux8fd9ux4e2aux5b8cux6574ux5bf9ux8c61ux63d0ux5347ux5230ux6807ux51c6-ell1}

这一步必须证明满射，不能只给稠密值域。

在 $c_0(I)$ 的单位球中，枚举所有有限支撑、有理坐标的向量： \[
q_1,q_2,\ldots .
\] 它们范数稠密，并包含单位坐标向量。

定义 \[
Q:\ell^1(\mathbb N)\to c_0(I),
\qquad
Qu=\sum_{n\ge1}u_nq_n.
\tag{21}
\] 级数绝对收敛，且 $\|Q\|=1$。

\subsection{\texorpdfstring{1. $Q$ 确实满射}{1. Q 确实满射}}\label{q-ux786eux5b9eux6ee1ux5c04}

给定 $x\in c_0(I)$，置 $r_0=x$。只要 $r_k\ne0$，选择 $q_{n_k}$ 满足 \[
\left\|q_{n_k}-\frac{r_k}{\|r_k\|_\infty}\right\|_\infty<\frac12,
\] 并置 \[
r_{k+1}=r_k-\|r_k\|_\infty q_{n_k}.
\]

于是 \[
\|r_{k+1}\|_\infty\le\frac12\|r_k\|_\infty,
\qquad
\sum_k\|r_k\|_\infty\le2\|x\|_\infty,
\] 且 \[
x=\sum_k\|r_k\|_\infty q_{n_k}.
\]

将重复索引的系数相加，就得到实际 $u\in\ell^1$，满足 \[
\boxed{
Qu=x,\qquad\|u\|_1\le2\|x\|_\infty.
}
\tag{22}
\]

因此 $Q$ 是满射。

\subsection{\texorpdfstring{2. 定义标准 $\ell^1$ 上的整张图}{2. 定义标准 \textbackslash ell\^{}1 上的整张图}}\label{ux5b9aux4e49ux6807ux51c6-ell1-ux4e0aux7684ux6574ux5f20ux56fe}

伴随明确为 \[
(Q^*a)_n=\langle q_n,a\rangle.
\]

定义 \[
\boxed{
\operatorname{gra}T
=
\{(u,Q^*a):a\in D,\ Qu=La\}.
}
\tag{23}
\]

非空性来自 $D\ne\varnothing$ 和 $Q$ 满射。单调性来自 \[
\langle u-v,Q^*(a-b)\rangle
=\langle La-Lb,a-b\rangle\ge0.
\]

\subsection{\texorpdfstring{3. 检查完整 $\ell^\infty$ 中的任意 polar 标签}{3. 检查完整 \textbackslash ell\^{}\textbackslash infty 中的任意 polar 标签}}\label{ux68c0ux67e5ux5b8cux6574-ellinfty-ux4e2dux7684ux4efbux610f-polar-ux6807ux7b7e}

设 \[
(v,v^*)\in\ell^1\times\ell^\infty
\] 与式（23）的全部行单调相关。

固定一条实际行 $(u_0,Q^*a_0)$。对所有 $k\in\ker Q$、$\theta\in\mathbb R$， \[
(u_0+\theta k,Q^*a_0)
\] 仍是实际行。测试全部正负 $\theta$，得到 \[
\langle k,v^*\rangle=0
\qquad(k\in\ker Q).
\tag{24}
\]

对 $x\in c_0(I)$，选任意 $u$ 满足 $Qu=x$，定义 \[
p(x)=\langle u,v^*\rangle.
\]

式（24）保证定义与原像选择无关，而且它是线性的。式（22）给出 \[
|p(x)|\le2\|v^*\|_\infty\|x\|_\infty.
\] 所以 \[
p\in(c_0(I))^*=\ell^1(I),
\qquad v^*=Q^*p.
\tag{25}
\]

现在对任意 $a\in D$，取 $u$ 满足 $Qu=La$，有 \[
0\le\langle v-u,v^*-Q^*a\rangle
=\langle Qv-La,p-a\rangle.
\]

由 $A$ 已证明的极大性， \[
p\in D,\qquad Qv=Lp.
\] 所以 $(v,v^*)$ 属于式（23）。

因此， \[
\boxed{T\text{ 在完整 }\ell^1\times\ell^\infty\text{ 中极大单调}.}
\tag{26}
\]

\subsection{4. 原命题的全部条件}\label{ux539fux547dux9898ux7684ux5168ux90e8ux6761ux4ef6}

若 $0\in\operatorname{dom}T$，则某个 $a\in D$ 满足 \[
La=Q0=0,
\] 与式（19）矛盾。因此 \[
T(0)=\varnothing.
\]

对任意 $(u,Q^*a)\in\operatorname{gra}T$， \[
\|u\|_1\ge\|Qu\|_\infty=\|La\|_\infty>\frac12,
\] 并且 \[
\langle u,Q^*a\rangle
=\langle La,a\rangle
>-\sigma
\ge-2\sigma\|u\|_1.
\]

故 \[
\boxed{
T(0)=\varnothing,\qquad
\mathcal V(T)\le2\sigma<\frac14.
}
\tag{27}
\]

\textbf{这里已经回到标准 $\ell^1$ 范数，没有使用等价范数，也没有只在某个较小对偶空间中证明极大性。}

\section{六、直接得到原始 CQ 下的求和反例}\label{ux516dux76f4ux63a5ux5f97ux5230ux539fux59cb-cq-ux4e0bux7684ux6c42ux548cux53cdux4f8b}

令 \[
f=Q^*g\in\ell^\infty,
\qquad
Bu=\langle u,f\rangle f.
\tag{28}
\]

因为 $Q$ 满射、$g\ne0$，所以 $f\ne0$。$B$ 全定义、非零、有界线性、正且秩一。

它的极大性可直接证明：若 $(z,z^*)$ 与整个 $B$ 图单调相关，测试 $(z+tw,B(z+tw))$，得到 \[
-t\langle w,z^*-Bz\rangle
+t^2\langle w,Bw\rangle\ge0
\qquad(t\in\mathbb R).
\] 让 $t$ 从正负两侧趋零，便有 \[
z^*=Bz.
\]

现在对任意 $(u,Q^*a)\in\operatorname{gra}T$，令 $t=r(a)$，则 \[
\langle u,f\rangle
=\langle Qu,g\rangle
=\langle La,g\rangle=t.
\]

所以，对\textbf{整个和图的全部实际行}， \[
\begin{aligned}
\langle u,Q^*a+Bu\rangle
&=\langle La,a\rangle+t^2\\
&=2t^2-1-\|\omega_{|t|-1}\|_2^2\\
&>1-\sigma>\frac78.
\end{aligned}
\tag{29}
\]

因此 \[
(0,0)\in(\operatorname{gra}(T+B))^\mu.
\] 但 $T(0)=\varnothing$，故 \[
(0,0)\notin\operatorname{gra}(T+B).
\]

同时， \[
\operatorname{dom}T\cap\operatorname{int}\operatorname{dom}B
=\operatorname{dom}T\ne\varnothing.
\]

于是得到完整的五项核验： \[
\boxed{
\begin{gathered}
T,\ B\text{ 均极大单调},\\
\operatorname{dom}T\cap\operatorname{int}\operatorname{dom}B\ne\varnothing,\\
(0,0)\in(\operatorname{gra}(T+B))^\mu,\\
(0,0)\notin\operatorname{gra}(T+B).
\end{gathered}}
\tag{30}
\]

由于 $T+B$ 单调，加入 $(0,0)$ 后得到其真单调扩张，故 $T+B$ 非极大单调。

由前述 $\Phi$ 和 $\Delta$ 的定义，还可以计算 \[
\boxed{\Delta_{T,B}(0,0)=\sigma\in(0,\infty).}
\] 理由是 \[
\Phi_B(0,p)=
\begin{cases}
\frac{\lambda^2}{4},&p=\lambda f,\\
+\infty,&p\notin\mathbb Rf,
\end{cases}
\] 当 $p=\lambda f$ 时，该公式来自 $\sup_{s\in\mathbb R}(\lambda s-s^2)=\frac{\lambda^2}{4}$。当 $p\notin\mathbb Rf$ 时，存在 $v\in\ker f$ 使 $\langle v,p\rangle\ne0$，沿 $v$ 的实数倍取上确界即得 $+\infty$。另一方面，由 $\langle u,f\rangle=r(a)$ 及 $Q$ 满射，
\[
\Phi_T(0,-\lambda f)
=\sup_{|t|>1}\bigl(1+\|\omega_{|t|-1}\|_2^2-t^2-\lambda t\bigr).
\]
实际参数 $t\downarrow1$ 和 $t\uparrow-1$ 给出 \[
\Phi_T(0,-\lambda f)\ge\sigma+|\lambda|,
\qquad
\Phi_T(0,0)=\sigma.
\] 因而
\[
\Phi_T(0,-\lambda f)+\Phi_B(0,\lambda f)
\ge\sigma+|\lambda|+\frac{\lambda^2}{4},
\]
且 $\lambda=0$ 时取等号，所以 $\Delta_{T,B}(0,0)=\sigma$。

\end{document}
